\section{Concept}\label{concept}

\textbf{Product type.}
\textit{Distributed Sensor Hub} is a decentralized peer-to-peer software platform for distributed sensor monitoring. It combines:
\begin{itemize}
  \item a set of autonomous, containerizable peer nodes that collect local readings, synchronize state in a distributed cluster, and replicate data without a central coordinator;
  \item an HTTP API for observability and integration;
  \item a read-only web UI for monitoring.
\end{itemize}


\textbf{Use case description.}
\textbf{What is the software doing?}
Each node maintains active connections to its local sensors, continuously collects readings, and updates its own local state. In parallel, the node shares its sensor-state updates with peer nodes and receives their updates through a gossip-based push-pull protocol. As a result, every node progressively learns and replicates the state produced by other nodes, building a distributed shared view. Eventual consistency is guaranteed through timestamp-based conflict resolution, while the system remains operational under failures, delays, and temporary partitions.

\textbf{Where are the users?}
In a smart-city setting, users are in two main groups. Technical operators/developers manage container deployment and node setup, and they extend the system because sensor protocols are injectable and node-to-node connection policies are pluggable; they may also build tools that consume exported data. Information-oriented users access the resulting data views to monitor conditions and support practical decisions or downstream services.

\textbf{When and how frequently do they interact with the system?}
Interaction is role-dependent. Technical operators/developers interact mainly during setup, topology changes, protocol/policy development, testing, and incident handling; this is recurring but not continuous. Information-oriented users interact when they need situational updates, monitoring views, or derived outputs for operational decisions. In contrast, the distributed core operates continuously and autonomously: sensors generate readings periodically and nodes exchange synchronization updates continuously.

\textbf{How do they interact with the system?}
Technical operators/developers interact through deployment/runtime tooling (local process management or Docker orchestration), protocol/policy implementation components, and programmatic inspection via HTTP endpoints. Information-oriented users interact through read-only monitoring views, typically via the web UI or external dashboards/tools that consume exported data.

\textbf{Does the system need to store user data? Which? Where?}
No user-specific data is stored. The system handles technical runtime data only: timestamped sensor readings, replicated node state, synchronization metadata, and membership/peer information. Data is maintained locally by each node (in-memory in the current scope) and propagated across peers by the gossip protocol.

\textbf{Multiple roles}
The model includes three technical roles:
\begin{itemize}
  \item \emph{Node (core role)}: ingests local sensor data, maintains local state, and participates in peer synchronization and replication;
  \item \emph{Technical operator/developer}: deploys and configures nodes, develops injectable sensor protocols and pluggable connection policies, and validates runtime behavior;
  \item \emph{Information-oriented user/client}: consumes read-only state views for monitoring and downstream operational use, without participating in replication.
\end{itemize}

\textbf{Why distribution is needed.}
Distribution is a core requirement, not an implementation detail. In a smart-city environment, sensing points are geographically distributed across different physical areas, so data is naturally produced at multiple nodes rather than at a single central site. Each node contributes local observations and shares them with peers, allowing the system to build a replicated global view without relying on a permanent coordinator.

Distribution is also required for fault tolerance. Nodes share sensor state, membership knowledge, topology information, and replication metadata through peer-to-peer communication, while remaining able to operate autonomously if other peers become unreachable. This allows the system to continue local ingestion and observation during node crashes, delays, or temporary network partitions, and to converge again when communication is restored.

Finally, distribution is essential because the project is meant to expose real distributed-systems behavior, including delay, partial failure, recovery, and eventual consistency. These aspects are part of the project goals and must be observable rather than hidden behind a centralized architecture.

The detailed behavioral requirements, protocol constraints, and acceptance criteria are specified in \cref{requirements}.

\clearpage
